%% Responsable: Sebastián López
\section{Algoritmo exacto}
\label{sec:exacto}

El exacto es un backtracking con cota (D3, D5). Está en el módulo
\texttt{backtracking.py} de \texttt{busqueda\_prompt/algoritmos}, se comprobó contra la fuerza bruta y
se midió con el modelo sintético (D7). Con el modelo real no se ejecutó: esta sección no
reporta ninguna precisión real ni afirma un óptimo global o de un subcatálogo.

\subsection{Descripción}

El árbol de búsqueda tiene un nivel por ranura: un nodo de nivel $\ell$ es una asignación
parcial de las ranuras $1,\dots,\ell$ y sus hijos prueban las $k_{\ell+1}$ opciones de la
siguiente, en orden creciente. Las hojas son las configuraciones de $\conf$, recorridas en
orden lexicográfico, el mismo de la fuerza bruta. En cada hoja la configuración se evalúa
instancia por instancia, y se lleva la mejor cantidad de aciertos conocida, $a^*$.

\begin{algorithm}[h]
\caption{Backtracking con cota}
\label{alg:backtracking}
\begin{algorithmic}[1]
\Require ranuras $1,\dots,L$ con $k_i$ opciones, lote $I=(x_1,\dots,x_n)$, cota parcial $U$ opcional
\State $a^*\gets-1$;\quad $c^*\gets$ ninguna
\Procedure{Explorar}{$\ell$, $p$} \Comment{$p$ asigna las ranuras $1,\dots,\ell$}
  \If{$\ell=L$}
    \State $a\gets 0$
    \For{$j=1,\dots,n$}
      \If{$a+(n-j+1)\le a^*$} \Return \Comment{cota D5: no puede superar a $c^*$}
      \EndIf
      \If{$p$ cumple $x_j$} $a\gets a+1$ \Comment{una generación}
      \EndIf
    \EndFor
    \If{$a>a^*$} $a^*\gets a$;\quad $c^*\gets p$
    \EndIf
    \State \Return
  \EndIf
  \If{$U$ dada, $\ell>0$, $a^*\ge0$ y $U(p)\le a^*$} \Return \Comment{poda por estructura (D6)}
  \EndIf
  \For{$o=0,\dots,k_{\ell+1}-1$}
    \State \Call{Explorar}{$\ell+1$, $p\cup\{\ell+1\mapsto o\}$}
  \EndFor
\EndProcedure
\State \Call{Explorar}{$0$, $\varnothing$};\quad \Return $c^*$, $a^*/n$
\end{algorithmic}
\end{algorithm}

La cota D5 no supone nada del oráculo: una instancia aún no evaluada aporta a lo sumo un
acierto, así que $a+(n-j+1)$ acota los aciertos finales de la configuración. Vale con el
modelo real y con el sintético, con o sin ruido.

La cota parcial $U$ solo se usa con el modelo sintético, donde la estructura de reglas es
conocida (D6). Para una asignación parcial $p$, las reglas alcanzables son las de las
opciones elegidas más todas las de las ranuras libres; $U(p)$ cuenta las instancias cuyas
reglas exigidas están todas en ese conjunto. Ninguna configuración que complete $p$ cubre una
regla fuera de él, y los conflictos solo quitan reglas, así que $U(p)$ es una cota superior.
Con ruido deja de valer y el código la rechaza. Con el modelo real no hay un $\sigma$ conocido
y no se poda con supuestos.

\subsection{Correctitud}

\begin{teorema}[Correctitud]\label{teo:exacto-correcto}
Si $U$ es una cota superior válida, o no se usa, el algoritmo~\ref{alg:backtracking}
devuelve la primera configuración del orden lexicográfico con la mayor precisión, la misma
que la fuerza bruta.
\end{teorema}
\begin{proof}
Sea $c^\circ$ la primera configuración del orden con el máximo de aciertos $a^\circ$. Toda
configuración anterior tiene menos de $a^\circ$ aciertos, de modo que al llegar a $c^\circ$
se cumple $a^*<a^\circ$. La cota D5 no la abandona: en cada paso $j$, sus aciertos finales
cumplen $a^\circ\le a+(n-j+1)$, y por tanto $a+(n-j+1)>a^*$. Ningún ancestro de $c^\circ$
se poda por $U$: si $p$ es ancestro, $U(p)\ge a^\circ>a^*$. Entonces $c^\circ$ se evalúa
completa y pasa a ser $c^*$. Después, ninguna configuración tiene más de $a^\circ$ aciertos y
el reemplazo exige $a>a^*$, así que $c^*$ ya no cambia. Las configuraciones que se
abandonan o se podan no podían superar estrictamente a $a^*$, así que no se pierde ninguna
que la fuerza bruta habría elegido.
\end{proof}

Las pruebas de \texttt{tests/test\_backtracking.py} comprueban el teorema: en 70 instancias
sintéticas, con conflictos, con ruido y con la cota estructural, la ganadora y la precisión
coinciden con las de la fuerza bruta. También comprueban, en 10 casos sintéticos, que $U$ nunca
queda por debajo de los aciertos reales en ninguna asignación parcial.

\subsection{Costo $T(n)$}

El costo se mide en generaciones, como en la sección~\ref{sec:aproximados}: cada vez que se
pregunta si $p$ cumple $x_j$ se hace una generación (D11).

\begin{teorema}[Costo]\label{teo:exacto-costo}
El algoritmo~\ref{alg:backtracking} hace a lo sumo $|\conf|\cdot n$ generaciones, y esa cota
se alcanza. Visita a lo sumo $\sum_{\ell=0}^{L}\prod_{i\le\ell}k_i$ nodos, que es menos de
$2|\conf|$ si todo $k_i\ge2$. Sin $U$, el tiempo total es $O(|\conf|\cdot(n+L))$, y el
número de generaciones es $\Theta(|\conf|\cdot n)$ en el peor caso.
\end{teorema}
\begin{proof}
Cada hoja evalúa a lo sumo sus $n$ instancias y hay $|\conf|$ hojas. Para el número de nodos,
el nivel $\ell$ tiene $\prod_{i\le\ell}k_i\le|\conf|/2^{L-\ell}$ nodos si todo $k_i\ge2$, y la
suma geométrica queda por debajo de $2|\conf|$. Un nodo interno cuesta $O(1)$ sin $U$, y una
hoja $O(n+L)$, contando la copia de la configuración, así que el tiempo es
$O(|\conf|\cdot(n+L))$. La cota se alcanza con un oráculo en el
que ninguna configuración cumple ninguna instancia: $a^*$ queda en $0$ tras la primera hoja,
y en cada hoja la condición de abandono $0+(n-j+1)\le0$ nunca se cumple, así que todas las
hojas evalúan sus $n$ instancias.
\end{proof}

En el mejor caso, si la primera configuración cumple todo, $a^*=n$ y las demás se abandonan
antes de su primera generación: el costo es $n$. Ambos extremos están cubiertos por pruebas
automáticas. Con $U$, cada nodo interno suma $O(n\cdot|R|)$ operaciones sin generaciones.

Con $L$ y los $k_i$ como parte de la entrada, $|\conf|=\prod_ik_i\ge2^L$ crece de forma
exponencial, lo cual es coherente con la NP-dureza de la sección~\ref{sec:reduccion}: la poda
mejora casos concretos, pero no cambia el peor caso. En el caso del proyecto, con $|\conf|$
fijo en $32\,768$, el costo es lineal en $n$ con una constante enorme.

\subsection{Medición con el modelo sintético}

El barrido sigue D10: $|\conf|=256$ fijo (4 ranuras de 4 opciones) con
$n\in\{10,25,50,75,100\}$, y $n=50$ fijo con $|\conf|=2^L$ para $L$ de 4 a 10. Cada punto
promedia 5 semillas de instancias con 12 reglas, 3 reglas por instancia, densidad $0{,}2$ y 3
conflictos. En todas las semillas el exacto coincidió con la fuerza bruta. Se ejecuta con
\texttt{python -m busqueda\_prompt.experimentos.medir}, que escribe
\texttt{resultados/exacto/barrido\_n.csv}, \texttt{barrido\_c.csv} y \texttt{ajuste.json}.

\begin{center}
\begin{tikzpicture}
\begin{loglogaxis}[
  width=0.48\textwidth, xlabel={$n$}, ylabel={generaciones},
  title={$|\conf|=256$ fijo}, legend pos=north west, legend style={font=\scriptsize},
  log basis x=10, log basis y=10]
\addplot+[mark=square*] table[x=n, y=fuerza_bruta, col sep=comma]{../resultados/exacto/barrido_n.csv};
\addplot+[mark=*] table[x=n, y=backtracking, col sep=comma]{../resultados/exacto/barrido_n.csv};
\addplot+[mark=triangle*] table[x=n, y=backtracking_estructural, col sep=comma]{../resultados/exacto/barrido_n.csv};
\legend{fuerza bruta, cota D5, D5 y estructural}
\end{loglogaxis}
\end{tikzpicture}\hfill
\begin{tikzpicture}
\begin{loglogaxis}[
  width=0.48\textwidth, xlabel={$|\conf|$}, ylabel={generaciones},
  title={$n=50$ fijo}, legend pos=north west, legend style={font=\scriptsize},
  log basis x=2, log basis y=10]
\addplot+[mark=square*] table[x=C, y=fuerza_bruta, col sep=comma]{../resultados/exacto/barrido_c.csv};
\addplot+[mark=*] table[x=C, y=backtracking, col sep=comma]{../resultados/exacto/barrido_c.csv};
\addplot+[mark=triangle*] table[x=C, y=backtracking_estructural, col sep=comma]{../resultados/exacto/barrido_c.csv};
\legend{fuerza bruta, cota D5, D5 y estructural}
\end{loglogaxis}
\end{tikzpicture}
\end{center}

\begin{center}
\begin{tabular}{lcc}
\toprule
Pendiente log-log de las generaciones & frente a $n$ & frente a $|\conf|$ \\
\midrule
Fuerza bruta y peor caso & $1{,}00$ & $1{,}00$ \\
Backtracking con cota D5 & $1{,}03$ & $0{,}52$ \\
Backtracking con D5 y cota estructural & $1{,}19$ & $0{,}28$ \\
\bottomrule
\end{tabular}
\end{center}

\paragraph{Lectura.} Frente a $n$, la cota D5 ahorra una fracción casi constante, entre el
45 y el 49\,\% con $|\conf|=256$ (14\,048 frente a 25\,600 generaciones con $n=100$), y el costo sigue
siendo lineal en $n$, como predice el teorema. Frente a $|\conf|$, el ahorro crece: con
$|\conf|=1\,024$ la cota D5 gastó 7\,195 generaciones frente a 51\,200, porque al encontrar
pronto una configuración con muchos aciertos casi todas las demás se abandonan tras pocas
instancias. La cota estructural poda subárboles enteros y baja más el gasto, aunque su
pendiente frente a $n$ es algo mayor porque su ventaja relativa es mayor con $n$ pequeño.

Estas pendientes describen este generador y estos parámetros; no son una garantía. Al crecer
$L$ con densidad fija, cada instancia se cubre con más facilidad y la precisión máxima sube,
lo que favorece la poda, y la serie frente a $|\conf|$ no es monótona (512 gastó menos que
256). El peor caso, que se midió con el oráculo sin aciertos, mantiene la pendiente $1$ en
los dos ejes. Los tiempos de pared fueron de $0{,}004$ a $0{,}09$ s por ejecución y siguen a
las generaciones; con el modelo real dominan las generaciones.

\subsection{Alcance con el modelo real}

Con $1{,}99$ s por generación medidos en T4, el peor caso del espacio completo serían
$3\,276\,800$ generaciones, unos 75 días de GPU. En un subcatálogo con las 3 mejores opciones
por ranura (243 configuraciones, D8) el peor caso serían $24\,300$ generaciones, unas 13 h, y
con 2 por ranura (32 configuraciones), $3\,200$, cerca de 1,8 h. Ninguno se ejecutó en esta
entrega, porque el presupuesto $B=14\,400$ se dedicó a la búsqueda por coordenadas (D17). Un
óptimo de subcatálogo solo sería óptimo en ese subcatálogo, y el subcatálogo mismo saldría de
la búsqueda, no de una demostración.
